\documentclass[openany,fontset=fandol]{ctexbook}
\usepackage{amsmath,amssymb}
\usepackage[unicode,hidelinks]{hyperref}
\newcommand{\note}[1]{\footnote{译注：#1}}
\begin{document}
\title{代数几何基础第五卷预备笔记\\引理13.3：一般截面的归纳}
\author{据用户上传英译编辑稿译出}
\date{2026年11月17日}
\maketitle
\chapter*{范围与校订状态}
本批连续译出底本第51页“This is a regret to No. 9”起至第52页“of a presentation (enforme)”止，共两个原段。下一处为第52页“It remains to prove the corollary announced of 8.8 in the case $m=1$”。

本批给出原稿所用超平面截面断言及向一般余维的归纳说明；所依赖的$m=1$情形尚待下段证明，不能据本批称引理13.3或定理13.2已证。8.8缺文与原稿记号疑点仍保留；正文仅整理可由紧邻上下文确定的错置，并以脚注保存原文及依据。未并入原合订本。
\chapter*{§13 初等态射与M. Artin定理（续）}
这是对第9节的一项补记，而它本身又来自第8节后面的补记。\note{底本在这些位置使用regret、regrets，按上下文是希望补入先前章节的结果，故译为“补记”，并非已正式编号的勘误。}
沿用8.8的记号，并假设$X$不可约；对于我们正在讨论的$p^0$问题（？），这样做是允许的。\note{“$p^0$-problem (?)”为底本原有的难辨记号与疑问，未擅自还原成某一性质。本批所用底本在8.7后没有完整的8.8条目，故此引用并不表示8.8缺文已经恢复。}
如果
\[
\operatorname{codim}T\geq2,
\]
或者$h^{\mathrm{sing}}\to T$一般平展，那么对于经过$x$的一般超平面$H_\eta$，截面$X^{(1)}_\eta$是光滑的，维数为$n-1$；这里须除去$f(X)=\{x\}$的情形（于是$n=0$）。\note{底本此处印$h^{\mathrm{sing}}\to T$，第52页后续证明改印$X^{\mathrm{sing}}\to T$，本批不将两种记号静默统一。原句“the hyperplane $H_\eta$ passing through a generic $x$”词序含混；上一段及下一页均把$x$固定、令$\eta$为经过$x$的超平面参数空间的一般点，正文据此整理为“经过$x$的一般超平面”，不另要求$x$是$X$或$P$的一般点。末句原印$f(x)=\{x\}$，按其所说$X$零维的例外及下一页对所有$z\in X$的讨论，应为$f(X)=\{x\}$，故正文订正。维数$n$与下句替换指标$n$重名，仍照底本保留。}
一般平展这一条件在$k$的特征为零时自动满足；或者，在将$f$换成$\Phi_n f$、其中$n\geq2$之后满足，参见第9节。\note{底本把这一条件说明插在前句中，正文为避免长括号而移作紧随其后的句子。这里保留其断言与引用，不把第9节的适用假设或所引结论视为已独立验证；上一批原文的$Q_n$与此处$\Phi_n$之间的记号不一致仍待核。}
\newpage
暂且承认这一结果，它就解决了上述补记中$m=1$的特殊情形。对$m$作归纳即可得到一般$m$的情形：注意，允许作基变换后，$X^{(m)}_\eta$可以先取一个经过$x$、余维为$m-1$的$L'$，再取由$H$给出的线性截面而得到；$L'$和$H$对于这些条件是相互独立的一般选择。\note{底本原句在“taking an $L'$”之后把后续动作放进长括号，语法不完整；这里仅按它明确列出的两步整理语序。所取的是相应关联参数空间上的一般点，而非任意两个线性子空间。没有由此声称一般纤维无需任何基变换即可相等。}
换成通常的严谨说法，就是把自己置于参数化二元组$(L',H)$的概形的一般点处，并取由$H$给出的线性截面；归纳假设提供了所需的光滑性。\note{底本括注末尾为“by taking the linear sections by $H$ which is smooth by inductive assumption”，其中which的指代不清。归纳结构中，应先用归纳假设处理余维$m-1$的截面，再用刚承认的$m=1$情形处理最后一个超平面截取，不能仅凭归纳假设就断言最后的截面光滑。正文保留其归纳说明，脚注明确这一步仍依赖尚待证明的$m=1$结果。}

这种类型的论证此前已用来推广2.6，例如推广到任意余维$m$的线性截面。值得把它在一般情形下一次性清楚写出，以便以后引用，而不必每次都进入一份正式表述中那些有些繁重的细节［法语；此处有空缺］。\note{底本末尾为“a bit heavy [Fr] [\quad] of a presentation (enforme)”。保留其未辨空缺；enforme按“en forme”的语义译为正式表述。这一段是作者关于编写方式的说明，不意味着本批已经补成了一项一般定理及其证明。}
\end{document}
